%% Chapter 6 - Results and Validation
%% Planned sections (master structure): 6.1 Research results; 6.2 Product and
%% technical results; 6.3 User or stakeholder validation; 6.4 Integration
%% results; 6.5 Summary of findings.
%%
%% CANONICAL RESULTS SOURCE (author-confirmed, 2026-10-06): the 1500-record
%% evaluation in "Thesis/Reserch Paper/NeuroShell_IRE - Reserch paper.tex"
%% (N=1500, 600 well-formed + 900 flawed).
%%   Exp1 (raw vs guarded): Gemma 86/72/14; Qwen 82/64/15 (DAcc/Catch/FP).
%%   Exp2: same guarded figures head-to-head.
%%   Exp3 (ablation, both backends): validation stack -30/-24 catch,
%%   hallucination validator -9/-14 catch; several layers increase catch on
%%   removal while raising FP (scope, network, self-consistency, cache).
%%   Exp4 (per-class catch): 100/95.6/87.5/81.3/78.1/59.4/46.7/40.
%% All single-run point estimates, no confidence intervals. The 200-record run
%% at git HEAD was a SMOKE TEST only; the 1500-record evaluation is the
%% legitimate one. Latency (p50 16.5 s, p95 70.9 s, cache 44 ms) is only in the
%% final paper, not the 1500 draft -- reconcile with the author before quoting.
\chapter{Results and Validation}\label{ch:results}

This chapter reports the measured outcomes of the study. Every figure is a
point estimate computed from the committed experiment record by the
reproduction command of Section~\ref{sec:reproducibility}, using the metrics
of Section~\ref{sec:measures} and the procedure of Section~\ref{sec:procedure}.
All measurements are single-run: no confidence intervals are attached, and the
comparison rule of Section~\ref{sec:data-analysis} (a difference outside the
$\pm$1--2~pp recorded run-to-run noise band is read as a layer effect) applies
throughout.

\section{Research Results}\label{sec:results-research}

\subsection{Experiment 1: Pipeline Value (RQ1)}\label{sec:exp1}

Raw, single-shot model output performs no flaw classification, so it trivially
achieves a 0\% catch rate on all 900 flawed records, and it misclassifies the
intent of 22\% of the 600 well-formed records with Gemma and 31\% with Qwen.
Adding the guarded pipeline raises decision accuracy $D_{\mathrm{Acc}}$
(Equation~\ref{eq:metric-dacc}) from 78\% to 86\% for Gemma and from 69\% to
82\% for Qwen, raises the flaw-catch rate (Equation~\ref{eq:metric-catch}) from
0\% to 72\% for Gemma and to 64\% for Qwen, and reduces the false-positive rate
(Equation~\ref{eq:metric-fp}) from 32\% to 14\% for Gemma and from 31\% to 15\%
for Qwen (Table~\ref{tab:exp1}). The largest improvement is in the catch rate,
by construction: raw output is never flagged at all.

\begin{table}[h]
\centering
\small
\caption{Experiment 1: raw model output against the guarded pipeline on the
full 1500-record set ($N{=}1500$, $D^+{=}600$, $D^-{=}900$).}\label{tab:exp1}
\begin{tabular}{@{}lccc@{}}
\toprule
Configuration & $D_{\mathrm{Acc}}$\% & Catch\% & FP\%\\
\midrule
Gemma raw      & 78 & 0  & 32\\
Gemma guarded  & 86 & 72 & 14\\
Qwen raw       & 69 & 0  & 31\\
Qwen guarded   & 82 & 64 & 15\\
\bottomrule
\end{tabular}
\end{table}

The result answers RQ1 directly: the guardrail stack converts output that
carries no machine-checkable safety property into output whose acceptance is
conditional on a verified conjunction of predicates (Section~\ref{sec:research-design}).
The catch-rate gap is categorical rather than incremental, because raw output
has zero catch by construction.

\subsection{Experiment 2: Backend Head-to-Head (RQ3)}\label{sec:exp2}

Under the identical guardrail stack, Gemma leads Qwen on decision accuracy
(86\% against 82\%) and on catch rate (72\% against 64\%) with comparable
false-positive rates (14\% against 15\%, Table~\ref{tab:exp2}). The guarded
architecture narrows the backend gap on accuracy from 9 percentage points
(78\% against 69\% raw) to 4 points, and holds the catch-gap at 8 points.

\begin{table}[h]
\centering
\small
\caption{Experiment 2: identical guardrail stack, both backends.}\label{tab:exp2}
\begin{tabular}{@{}lccc@{}}
\toprule
Backend & $D_{\mathrm{Acc}}$\% & Catch\% & FP\%\\
\midrule
Gemma (guarded) & 86 & 72 & 14\\
Qwen (guarded)  & 82 & 64 & 15\\
\bottomrule
\end{tabular}
\end{table}

\subsection{Experiment 3: Layer Ablation (RQ2 and RQ3)}\label{sec:exp3}

Table~\ref{tab:ablation} reports the eight configurations on the full 1500-record
set for both backends, with $\Delta$Catch relative to the fully guarded
baseline (Equation~\ref{eq:delta-catch}). Disabling the complete deterministic
validation stack (S5--S7) has the largest impact on catch rate: 30 points for
Gemma (72\% to 42\%) and 24 points for Qwen (64\% to 40\%). Among individually
ablated layers, the hallucination validator contributes most to catch rate
(9 points for Gemma, 14 points for Qwen). Its accuracy effect differs by
backend: largest for Qwen, whereas the scope guard has the largest accuracy
effect for Gemma.

Removing RBAC reduces both accuracy and catch in both backends while raising
FP, indicating that the role check improves accuracy and lowers false positives
simultaneously. The network validator is the second-largest accuracy
contributor (18 points Gemma, 24 points Qwen). The scope guard, network
validator, self-consistency gate and semantic cache each show a
\emph{rise} in catch on removal (2--10 points) together with a sharp rise in FP
(to 30--44\%). Per the reading rule of Section~\ref{sec:data-analysis}, these
positive deltas signal interactions rather than contributions: the layers act
as false-positive suppressors, and the pre-emptive-rejection hypothesis is
recorded as an open mechanism (Section~\ref{ch:discussion}).

\begin{table}
\centering
\small
\caption{Experiment 3: ablation on the full 1500-record set. $\Delta$Catch is
in percentage points relative to the fully guarded pipeline (a negative value
means the layer contributes); Noise is the approximate run-to-run variation in
percentage points. The validation-stack row is the joint S5--S7 ablation.}
\label{tab:ablation}
\setlength{\tabcolsep}{4pt}
\begin{tabular}{@{}lccccccccc@{}}
\toprule
 & \multicolumn{4}{c}{Gemma} & \multicolumn{4}{c}{Qwen} & \\
\cmidrule(lr){2-5}\cmidrule(lr){6-9}
Configuration & $D_{\mathrm{Acc}}$ & Catch & FP & $\Delta$Catch
& $D_{\mathrm{Acc}}$ & Catch & FP & $\Delta$Catch & Noise\\
\midrule
Full pipeline (baseline) & 86.0 & 72 & 14 & -- & 82.0 & 64 & 15 & -- & $\pm$1--2\\
No scope guard       & 65.0 & 74 & 42 & +2  & 68.0 & 68 & 30 & +4  & $\pm$1--2\\
No RBAC guard        & 75.0 & 65 & 20 & $-$7 & 57.0 & 57 & 17 & $-$7 & $\pm$1--2\\
No network validator & 68.0 & 80 & 40 & +8  & 58.0 & 72 & 32 & +8  & $\pm$1--2\\
No halluc. validator & 66.0 & 63 & 40 & $-$9 & 52.6 & 50 & 34 & $-$14 & $\pm$1--2\\
No validation stack  & 65.0 & 42 & 15 & $-$30 & 62.0 & 40 & 16 & $-$24 & $\pm$1--2\\
No self-consistency  & 68.0 & 78 & 42 & +6  & 58.0 & 74 & 36 & +10 & $\pm$1--2\\
No semantic cache    & 67.0 & 77 & 44 & +5  & 60.0 & 70 & 40 & +6  & $\pm$1--2\\
\bottomrule
\end{tabular}
\end{table}

\subsection{Experiment 4: Per-Class Catch (RQ2)}\label{sec:exp4}

On the dedicated deep-dive subset, catch is strongest exactly where a
deterministic rule exists (Table~\ref{tab:perclass}). Adversarial injection is
caught at 95.6\% by the lexical and structural layers of the detector, and RBAC
violations at 100\% by the deterministic pre-inference role check. Fabricated
parameters (87.5\%) and type mismatches (81.3\%) are detected by the S5--S6
syntactic rules, and fabricated CVEs (78.1\%) by the hallucination validator's
offline grounding. The residual weak classes are the semantic cases---contradictory
requests (59.4\%), out-of-scope requests (46.7\%) and ambiguous requests
(40.0\%)---where no deterministic signal exists and the pipeline relies on the
self-consistency gate. In the full set these three classes account for 339 of
the 900 flawed records (37.7\%). Per-class denominators therefore carry the
interpretation, as fixed in Section~\ref{sec:data-analysis}.

\begin{table}[h]
\centering
\small
\caption{Experiment 4: per-class catch rate and primary detection mechanism
(deep-dive subset, each class's own denominator).}\label{tab:perclass}
\begin{tabular}{@{}lcl@{}}
\toprule
Flaw class & Catch\% & Mechanism\\
\midrule
RBAC violation   & 100.0 & M2 role check\\
Adversarial inj. & 95.6  & M1 detector\\
Fabricated param.& 87.5  & S5--S6 rules\\
Type mismatch    & 81.3  & S5--S6 rules\\
Fabricated CVE   & 78.1  & hallucination validator\\
Contradictory    & 59.4  & self-consistency\\
Out-of-scope     & 46.7  & self-consistency\\
Ambiguous        & 40.0  & self-consistency\\
\bottomrule
\end{tabular}
\end{table}

\section{Product and Technical Results}\label{sec:product-results}

\subsection{Verification and testing}

The automated suite was re-derived by direct execution for this dissertation
(session 2026-10-07; Chapter~\ref{ch:implementation-testing}): 506 collected
tests, 506 passed, with 434 unit and 72 integration tests, and 76\% statement
coverage across \texttt{src/} (2421 statements, 582 missed; 17 modules at
100\%). The safety-critical core is the best-covered part of the system:
adversarial detection, intention schema, RBAC, scope guard and the regex and
network validators each sit at or above 93\%, while audit logging, semantic
cache and session context sit at 89--96\%. Static analysis
(\texttt{ruff}) reports 593 findings, all style and modernisation with no
correctness-level (E9) errors and no fatal findings. These results are the
product-level acceptance gate: they support the empirical measurements of
Section~\ref{sec:results-research} being representative of the committed system.

\subsection{Known limitation achieved}

The three semantic classes with the lowest catch in Experiment~4 (ambiguous,
out-of-scope, contradictory) are exactly the classes without a deterministic
signal; their residual misses constitute the principal functional gap of the
component and are analysed in Chapter~\ref{ch:discussion}.

\section{User or Stakeholder Validation}\label{sec:stakeholder-validation}

The study involves no human participants (Section~\ref{sec:ethics}), so no
usability or user-preference study was conducted. Stakeholder validation is
instead evidence-based, through two channels. First, the verification gate of
Section~\ref{sec:product-results}---a deterministic, live-free test suite that
must pass before any measurement is accepted---validates the component against
its stated requirements (Section~\ref{sec:requirements}) and the invariant of
Section~\ref{sec:research-design}. Second, the NeuroShell framework integration
review (Section~\ref{sec:integration-results}) validated the component's
contract against Component~2 of the framework on the receiving side.

\section{Integration Results}\label{sec:integration-results}

The C1-to-C2 integration analysis (session 2026-08-28; documented in the
project record) verified the contract hand-off against the real Component~2
engine at port 8002. The 9-value intent literal sets match exactly, and the
flat \texttt{PlannerContract} fields map 1:1 onto the Component~2 request. Live
checks confirmed: a well-formed in-scope contract returns a structured
planner output with an empty safety block; a rejected intent is refused by
Component~2 with its filter-stage error; an out-of-scope public target is
rejected by the Component~2 RFC-1918 gate; and a wrong API key is refused with
401. These observations confirm that the guarded IRE contract carries its
integrity through the downstream component, so the safety invariant of
Section~\ref{sec:research-design} remains meaningful at the framework boundary.
\TODO{The C1-to-C2 live verification used the available stand-in model at the
time; re-run the two-element happy-path check against the production
\emph{gemma4:e4b} backend once the model is provisioned on the evaluation
machine, and record the result in the committed experiment record.}

\section{Summary of Findings}\label{sec:summary-findings}

The measured outcomes answer the three research questions as follows. RQ1:
the guarded pipeline raises decision accuracy by 8--13 points, the catch rate
from 0\% to 64--72\%, and lowers the false-positive rate to 14--15\% over raw
output. RQ2: the deterministic validation stack is the dominant safety
contributor (24--30 catch points), the hallucination validator is the dominant
single layer (9--14 catch points), and the semantic classes remain the weakest
because they lack a deterministic signal. RQ3: the catch-rate ordering of
layers carries across both backends, but the accuracy contribution of
individual layers is backend-specific, so cross-backend generality claims
apply to catch ordering rather than to a single accuracy ranking. The
quantitative case is made by Table~\ref{tab:exp1}, Table~\ref{tab:exp2},
Table~\ref{tab:ablation} and Table~\ref{tab:perclass}, and is interpreted in
Chapter~\ref{ch:discussion}.